\documentclass{article}
\usepackage[preprint]{neurips_2025}
\usepackage[T1]{fontenc}
\usepackage[utf8]{inputenc}
\usepackage{microtype}
\usepackage{booktabs}
\usepackage{tabularx}
\usepackage{url}
\usepackage[hidelinks]{hyperref}
\usepackage{listings}
\makeatletter
\newenvironment{fixedtable}{\par\medskip\noindent\begin{minipage}{\linewidth}\def\@captype{table}}{\end{minipage}\par\medskip}
\makeatother
\lstset{basicstyle=\ttfamily\footnotesize,breaklines=true,columns=fullflexible,frame=single}
\newcommand{\TotalTokens}{3645}
\newcommand{\TotalSeconds}{37.971}

\newcommand{\ResultRows}{
1 & 465 & 264 & 15.366 & Saved \\
2 & 465 & 264 & 5.637 & Saved \\
3 & 465 & 264 & 5.670 & Saved \\
4 & 465 & 264 & 5.616 & Saved \\
5 & 465 & 264 & 5.682 & Saved \\
}
\newcommand{\CombinedRows}{
A & 1 & 465 & 264 & 15.366 & Saved & Met \\
A & 2 & 465 & 264 & 5.637 & Saved & Met \\
A & 3 & 465 & 264 & 5.670 & Saved & Met \\
A & 4 & 465 & 264 & 5.616 & Saved & Met \\
A & 5 & 465 & 264 & 5.682 & Saved & Met \\
B & 1 & 593 & 431 & 19.142 & Saved & Missed \\
B & 2 & 593 & 428 & 7.987 & Saved & Missed \\
B & 3 & 593 & 428 & 8.028 & Saved & Missed \\
B & 4 & 593 & 428 & 8.056 & Saved & Missed \\
B & 5 & 593 & 428 & 7.963 & Saved & Missed \\
}
\newcommand{\ExtendedFindings}{Batch B used 5,108 tokens and 51.176 seconds. All 5 attempts saved drafts and passed the code checks. The office-day check passed in 0/5 runs, and the route mailbox/case check passed in 0/5. The submission details appeared under participation, while response route stayed unknown. We count all five drafts as incomplete.

A likely cause is our own prompt. It says the official submission route may go under participation, it never defines participation as events such as an office day, and it does not mention office days. Our prompt allowed this. We have not changed the prompt or rerun, so we cannot say how much a clearer prompt would help. The saved Norwegian drafts also contradict themselves: the line on what a resident can do contains the mailbox and address, while the next line says the official route is not given. The office day does not appear at all.
}

\title{Hamar Planning Brief:\\An Evidence-Linked Walking Skeleton for Public Notices}
\author{\textbf{Kai Arne Hage}\\\textbf{Mohammad Almojtaba Elhajj A Elhajj}\\
Group 12, KIUA-2003\\University of Inland Norway}
\begin{document}
\raggedbottom
\maketitle
\begin{abstract}
Public planning notices ask residents to understand a proposal, identify its location,
and respond by a stated deadline. We propose a Norwegian draft-brief workflow that
extracts these facts with supporting passages and makes missing information visible.
The current Python skeleton loads a selected public-source extract, calls a local
model, formats a draft deterministically, and checks the output before saving it.
Five repeated runs on one short Ringgata extract passed the implemented checks.
A rejected model run and a synthetic injection case show why those checks cannot
replace human review. This report describes the working checkpoint, its limits,
and the plan for broader evaluation and stress testing before the examination.
\end{abstract}

\section{Problem, users, and definition of success}
The intended user is a Hamar resident seeking a traceable explanation of what a
planning notice proposes, where it applies, its stage, and when and how to respond.
The first example concerns Ringgata between Solvangvegen and Disenvegen
\citep{hamar}. The official local source and Norwegian output give the project a
concrete Innlandet focus. Resident need is a hypothesis: no user interviews or
usability measurements have been completed.

Scope is official Hamar-area land-use or transport planning notices inviting public
input. Grants, unrelated news, and completed decisions are out of scope. A reviewer
selects an extract before the workflow runs. The prototype combines structured
extraction, evidence, explicit uncertainty, and repeatable checks.

The first input contains three non-contiguous short passages selected and
source-checked by Codex on 28 October 2026, with no individual names or personal
contact details. Metadata records provenance and partial-source status. The response
route is omitted. Wider collection and redistribution require a source-policy review;
no general reuse license has been verified.

\textbf{Checkable success.} A run must save either an explicit failure report or a
Norwegian draft whose substantive fields have source evidence. For the selected
Ringgata extract, the draft should describe the cycleway and pavement proposal,
the two street endpoints, the start of planning work, and the stated input deadline
of 18 November 2026. Participation details and the response route must remain
\texttt{unknown} because they are absent from this input. The output must identify
the source, show a check date, and carry \texttt{needs\_human\_review}.
Mechanical checks establish format and locally checkable evidence. A person must
still assess meaning, completeness, and source freshness.

\section{Architecture and context as implemented}
Repository: \url{https://github.com/Kaiju77/hamar-planning-brief}.
It contains the context prompt,
Python code, synthetic fixtures, the selected real extract and metadata, saved runs,
tests, and reproduction instructions. The implementation uses Python's standard
library and Ollama's local generation endpoint \citep{ollama}.

\paragraph{Stage 1: load and bound the input.}
\texttt{run.py} reads the metadata JSON and selected text file. It validates the
check date and limits the source to 12,000 characters. The model receives only
\texttt{context/extract\_prompt.md} and the selected extract inside source
delimiters. It receives neither the full website nor conversation history, other
notices, credentials, or filesystem access. The URL is provenance, not a page
that the model independently fetches. Source selection and privacy screening are
currently human responsibilities.

\paragraph{Stage 2: extract structured facts.}
One local model call proposes a Boolean scope decision and six objects:
\texttt{proposal}, \texttt{location}, \texttt{planning\_stage}, \texttt{deadline},
\texttt{participation}, and \texttt{response\_route}. Each contains a string value
and an evidence passage. Instructions require unknown or ambiguous values instead
of guesses. A deadline may be normalized to an ISO date only when day, month, and
year are stated. The response route must be explicitly present, rather than
inferred from the website domain. The selected model is
\texttt{qwen2.5-coder:7b}, at temperature zero, in JSON output mode.

\paragraph{Stage 3: format a Norwegian draft.}
A deterministic template renders the facts. It leads with the proposal, location,
deadline and status, participation, and response route, followed by planning stage,
evidence, and source. There is no second drafting model. This departs from the
initial two-model proposal and removes one opportunity for new unsupported prose.
It also limits the quality of plain-language rewriting. Missing routes direct the
reader to the original notice, and partial-source status is visible.

\paragraph{Stage 4: check and save.}
Before any draft is saved, \texttt{checks.py} checks types, provenance format,
field evidence, and deadline validity. Evidence matching ignores case and repeated
whitespace. The normalized deadline must correspond to a recognized Norwegian
date string inside its own evidence. A code-computed status compares the deadline
with the recorded Europe/Oslo check date. A date-only deadline remains open through
that date; an unknown or ambiguous deadline has unknown status. Notices with precise
cutoff times require future handling. Failed checks save structured errors rather
than a draft. Exact passage matching does not establish semantic entailment.

\paragraph{Automation and execution limits.}
One command runs the stages without manual handoffs. Python reads and writes local
files, invokes the local HTTP service, and performs checks. The model has no tools
that execute commands or publish results. Each invocation gets a distinct directory
with raw response, counts, elapsed time, facts, and an input/prompt snapshot. There
is one model call, no automatic retry, an 800-token generation cap, and a 120-second
HTTP timeout. These are bounded settings, not yet a measured total token budget.
Ctrl+C stops the client; stopping the model may additionally require
\texttt{ollama stop}. Automatic publication is excluded and human approval is required.

\section{Walking-skeleton evidence and limitations}
The recorded final batch uses the response-route prompt and the selected Ringgata
extract, with check date 28 October 2026. Each of five Qwen runs passed mechanical
checks and saved a draft. The table is generated from the preserved
\texttt{run.json} files. Token counts come from Ollama; seconds are client wall time
through extraction and checking, excluding final file writes.

\begin{fixedtable}
\centering
\begin{tabular}{rrrrl}
\toprule
Run & Input tokens & Output tokens & Seconds & Outcome\\
\midrule
\ResultRows
\bottomrule
\end{tabular}
\caption{Five repetitions of one short real-source input, not five distinct notices.}
\end{fixedtable}

All five drafts contain the intended normalized deadline. Narrow reference checks
also find both location endpoints and the expected proposal and stage terms in
5/5 runs; participation and response route remain unknown in 5/5. These are
word-based proxies, not independent semantic grading. Two fields are only tested
as unknown on real input. Identical extracted facts in all five runs provide little
evidence about variance. Assistant inspection agrees
with the selected extract, but student review remains to be performed. The batch
used \TotalTokens{} input-plus-output tokens and \TotalSeconds{} seconds in total.
Longer first-run latency includes model loading, so this is not a controlled
performance comparison. A temperature of zero does not guarantee reproducibility
across model versions, machines, or runtime changes.

Earlier exploratory evidence is retained separately. One \texttt{llama3.2:3b}
response rejected the real notice as out of scope and produced invented quotations
and a non-ISO date. The pipeline saved a failure report. Because scope rejection
short-circuits later checks, this does not demonstrate detection of every invented
quotation. Five earlier Qwen successes predate the response-route field and are not
combined with the final batch as independent evidence.

Twelve unit tests pass. Seven original fixture cases exercise normal output,
missing and ambiguous dates, scope rejection, altered evidence, and two injection
outcomes. Four added tests cover malformed deadlines, date boundaries, missing
input, and preservation of repeated runs. A fifth robustness test mocks an
unavailable Ollama service and verifies failure logging without a draft.
The fixture backend supplies canned model
outputs: a test labeled as model resistance proves only that the checker accepts
that prewritten response. It is not an observed attack against a live model.

The known injection gap is explicit. A synthetic source contains an instruction
and an injected date. If a model adopts that date and quotes it, exact matching can
pass. A word-pattern warning is raised, but the draft is still saved for review.
The fixture demonstrates a weakness of evidence matching, not a solved defense.
The 21-word extract makes extraction close to copying; it does not test full notices.
Non-date values are not checked for semantic support: a fabricated location or
personal contact detail paired with a real quote can pass. Other limits include
untested multi-deadline reasoning, no live collection, and no resident usability
study. The workflow extends a prompt through provenance, computed status,
validation and failure reports, meeting the technical minimum of the brief
\citep{brief}; the evaluation plan must broaden before the examination.

\section{Plan, risks, and progression to the examination}
Our Compulsory 2 deadline is 7 November 2026 at 23:59. We will confirm the assigned
demonstration slot in Canvas by 2 November and prepare for the earliest possible
session on 9 November. We plan backwards from weeks 46--47 and the white-paper deadline of 3 December
2026 at 11:59 \citep{brief}.

\begin{fixedtable}
\small
\begin{tabularx}{\linewidth}{@{}p{0.20\linewidth}X@{}}
\toprule
Target & Deliverable or decision gate\\
\midrule
By 2 Nov & Review source and output together; confirm roles, template and live-demo slot in Canvas.\\
By 5 Nov & Freeze baseline cases and evaluation; rehearse a normal run and an explicit failure for a possible 9 November demo.\\
6 Nov & Edit the report, add genuine meeting notes and reproduce from a clean clone.\\
7 Nov, 23:59 & Submit reviewed Compulsory 2 report and repository link through Canvas.\\
9--22 Nov & Attend the mandatory live demonstration in weeks 46--47; log observed failures and feedback.\\
By 25 Nov & Complete two fault classes, implement repairs and rerun all regression cases; revise architecture to match the built system.\\
26 Nov--2 Dec & Final quantitative and qualitative analysis, eight genuine meeting records, report editing and reproduction check.\\
3 Dec, 11:59 & Submit the 7--8-page white paper via Wiseflow; tag the submitted public repository commit \texttt{exam} and preserve it.\\
\bottomrule
\end{tabularx}
\caption{Proposed milestones; intermediate dates are group planning targets.}
\end{fixedtable}

\begin{fixedtable}
\small
\begin{tabularx}{\linewidth}{@{}p{0.31\linewidth}X@{}}
\toprule
Likely failure & Planned response\\
\midrule
Unsupported facts or weak Norwegian & Human comparison with the source; expand reference cases and reject unsupported claims.\\
Ambiguous or changed deadline & Preserve uncertainty and retrieval date; check the original before acting.\\
Ollama unavailable or invalid output & Save a failure report; diagnose the service and rerun explicitly. Use fixtures only as a labeled offline demonstration.\\
Injection or personal details in output & Retain human approval; screen inputs and outputs and test stronger controls. The current warning is insufficient.\\
Source reuse or access restrictions & Stop wider collection; use minimal attributed extracts and synthetic fixtures pending review.\\
\bottomrule
\end{tabularx}
\caption{Risk assessment: failure modes and responses.}
\end{fixedtable}

The planned stress classes are \textbf{input perturbation} and
\textbf{dependency failure}. Input tests will embed misleading instructions and
conflicting facts in synthetic notices. Dependency tests will simulate a refused
connection, timeout, or malformed model response. We will record the failure path,
repair the responsible stage, and demonstrate that the same case and the rest of
the regression set pass afterwards. These experiments are plans, not completed
security results.

The examination set will cover typical, missing, conflicting and out-of-scope cases.
Format and dates are checked by code; meaning and clarity require documented human
assessment. Further exam work includes explicit read/write permissions,
secret exclusion, an enforced total token budget, and a measurable halt policy.
A proposed halt gate is three consecutive invalid runs followed by operator review;
it is not implemented. Keep at least eight genuine meeting records. The early demo
will show the current skeleton and baseline limitations; later repairs and feedback
will be documented in the final paper.

\clearpage
\begin{thebibliography}{9}
\bibitem[Hamar kommune(2026)]{hamar}
Hamar kommune. Oppstart reguleringsplan fv. 187 Ringgata i Hamar.
Published/modified 26 October 2026; accessed 28 October 2026.
\url{https://www.hamar.kommune.no/lokaldemokrati/oppstart-reguleringsplan-fv-187-ringgata-i-hamar}.
\bibitem[University of Inland Norway(2026)]{brief}
University of Inland Norway. \emph{Agentic Workflows for Innlandet: Building and
Verifying Coding Agent Harnesses for Local Needs}. KIUA-2003 project brief,
Autumn 2026. Supplied course document.
\bibitem[Ollama(2026)]{ollama}
Ollama. Generate a response: local generation API documentation. Accessed
28 October 2026. \url{https://docs.ollama.com/api/generate}.
\end{thebibliography}

\appendix
\section{Agent prompt excerpts and corrections}
This appendix contains representative excerpts from the student's actual assistant
conversations. They document assistance, not completed human approval of the draft.
The excerpt section occupies one page, within the two-page prompt-log limit.

\paragraph{1. Choosing a feasible scope (27 October).}
\begin{quote}\small
\textbf{Kai:} ``Something that is relatively easy to pull off but still relevant and
good enough for a great exam/grade''\par
\textbf{Codex:} ``I would still choose the Hamar planning-consultation brief, with
one refinement: focus on helping a resident answer `What is being proposed, where,
and by when can I respond?' ''
\end{quote}\par
The excerpt records a proposed scope, not a demonstrated resident need.

\paragraph{2. Correcting the date-status design (27 October).}
\begin{quote}\small
\textbf{Claude review supplied by Kai:} ``Add a deadline\_status field (open,
closed, unknown) computed in code by comparing deadline with retrieved\_at.''\par
\textbf{Codex:} ``the program should compare the deadline with the current recorded
check date, not only the retrieval date.''
\end{quote}\par
The revised checker tests before, on and after the deadline. The source retrieval
date is retained separately.

\paragraph{3. Replacing an untested assumption with observed evidence (28 October).}
\begin{quote}\small
\textbf{Claude handoff supplied by Kai:} ``The Ollama backend is written but untested.''\par
\textbf{Codex:} ``The first real run was rejected: Llama marked the notice out of scope
and invented supporting quotations.''
\end{quote}\par
The failed response was retained. Subsequent Qwen runs are separate evidence rather
than a correction to the failed transcript. Excerpts reproduce words from the
conversation; Markdown emphasis and code formatting are typeset as ordinary text.

\paragraph{Assistance and editorial responsibility.}
Claude and Codex performed substantial coding, testing, execution and drafting.
Kai selected the direction, supplied critique and authorized preparation. The report
remains an editorial draft until both students review and edit it. No partner
approval, resident interview or completed team meeting is inferred.

\clearpage
\section{Working-session and meeting records}
These records describe traceable individual working sessions involving Kai and
assistants. They do not claim that Mohammad attended or that a formal group meeting
took place. Group meeting records should be added as meetings actually occur.

\subsection{27 October 2026: scope and source investigation}
\textbf{Participants evidenced:} Kai and Codex in the coursework conversation.
\textbf{Purpose:} interpret the new brief and choose a feasible Innlandet problem.
\textbf{Decisions:} investigate a Hamar planning brief, start from selected public
input, make output checkable, and retain human review. \textbf{Revisions:} a Claude
review supplied by Kai led to an explicit response route, current check date,
field-keyed evidence, a scope rule and a failure-report contract.
\textbf{Open items:} group arrangement, exact Canvas deadline, source reuse and
the supplied LaTeX template.

\subsection{28 October 2026: handoff, implementation, and evidence}
\textbf{Participants evidenced:} Kai and Codex, with a Claude handoff supplied by Kai.
\textbf{Purpose:} inspect the working skeleton and obtain real-run evidence.
\textbf{Actions:} reproduce seven tests; identify a malformed-deadline crash;
add four robustness tests; preserve unique outputs; verify a minimal source extract;
record one rejected Llama run, an earlier five-run Qwen batch, and a final five-run
Qwen batch after the output-contract change. \textbf{Decisions:} keep deterministic
drafting and document its departure from the initial plan; keep the injection gap
visible; publish a reproducible repository under Kai's authorization.
\textbf{New information:} Kai confirmed Group 12 membership and a Canvas deadline
of 7 November at 23:59. \textbf{Open actions:} student source/output review, partner
discussion and contribution confirmation, editorial revision, and checking the
Canvas-specific template. These actions are pending, not recorded as completed.

\subsection{Plan for genuine group records}
Agree the first group-meeting date; later tentative dates are 2, 5, 9, 12, 16, 23, and 30 November,
adjusting dates to availability. For each, record actual attendees, alternatives
considered, decisions, evidence reviewed, and next actions. These are proposed
dates only; they do not count as completed meeting reports.

\section{Reproduction and editorial status}
Install Python 3.10+ and Ollama, then pull \texttt{qwen2.5-coder:7b}. From the project
root, use:
\begin{lstlisting}
python run.py --meta examples/ringgata.meta.json --input examples/ringgata.txt --backend ollama --checked-at 2026-10-28
python -m unittest discover -s tests -v
python evaluate.py --backend fixture
python evaluate.py --backend ollama --count 5
python tools/build_results.py
\end{lstlisting}
The date above reproduces the recorded check; a current operational run should pass
the actual Oslo check date. Fixture mode is offline and uses canned extraction,
not a real model. \texttt{tools/build\_results.py} regenerates the report table from
the fixed saved final batch. Clean-clone reproduction is recorded separately in the
repository when completed.

The report uses the official NeurIPS 2025 style with the preprint option and no
NeurIPS checklist. No separate Canvas template has been supplied for comparison.
Students must inspect the source and output, confirm contributions and meeting
records, and edit the prose before submission. The local source check and generated
results are assistant work until independently reviewed. This status is preserved
to avoid claiming activities that have not yet occurred.
\end{document}
